% !TeX root = ../../Thesis.tex
\section{Introduction}
%credit: Credit: NASA/JPL/Space Science Institute, the Cassini imaging team home page, http://ciclops.org.
Their potential to host extraterrestrial life has made icy moons, 
 such as the Jovian moon Europa and Saturn's moon Enceladus, prime targets for
 exploration in the upcoming decades. The Cassini-Huygens mission 
 discovered active cryovolcanism on Enceladus in $2005$ in the form of
 water rich plumes erupting from the south polar region (~\fig{fig:eneladus_fountains}). These plumes 
 consist mainly of water vapour, ice particles along with dissolved gases with a source rate estimated to be
 $300\,$\si {kg/s} and supply most of the material 
 making up Saturn's E ring.\cite{fontesIcyPlumesEnceladus, noneExperimentalSimulationAssessment2022}

The mechanisms of cryovolcanism and the disparity between the 
 recorded flow behaviour and constituents of the plumes are subject
 to ongoing research. 
 Various models aim to explain the plume mechanics and the observed 
 flow behaviour - explosive boiling due to exposure of the subsurface 
 ocean to vacuum due to 
 cracks opened by tidal stresses; simulating the 
 fissures in the surface as convergent-divergent (de Laval) nozzles that
 accelerate the vapor to the observed supersonic speeds; or channels inside
 the icy crust acting as a de Laval nozzle where the grains are 
 initially decelerated 
 due to frequent wall collisions before being re-accelerated by the 
 vapour flow.~\cite{noneExperimentalSimulationAssessment2022,yeohUnderstandingPhysicsEnceladus2015}
 To be able to use plumes as a surrogate to learn about the subsurface 
 ocean, a better 
 understanding of their driving processes is key.   
\begin{figure}[h]
  \centering
  \caption{Enceladus and its cryovolcanic plumes. Credit: NASA/JPL/Space Science Institute.}
  \begin{subfigure}[b]{0.45\textwidth}
    \includegraphics[width=\linewidth]{./Chapters/Introduction/figs/enceladus_plumes}
    \caption{Distant view of Enceladus.}
    \label{fig:enceladus_plumes}
  \end{subfigure}
  \hfill
  \begin{subfigure}[b]{0.45\textwidth}
    \includegraphics[width=\linewidth]{./Chapters/Introduction/figs/enceladus_fountains}
    \caption{Enhanced, colour-coded image of the plumes outlining their extreme extents.}
    \label{fig:eneladus_fountains}
  \end{subfigure}
  
\end{figure}

Icy moon experiments, in an effort to contribute to the understanding 
 of the fluid dynamics of the supersonic plumes, an experimental setup is under
 development at the FH Aachen - University of Applied Sciences. In collaboration with the 
 Rocket Experiments for University Students (REXUS) programme of the 
 German Aerospace Center (DLR) and the Swedish Space Agency (SNSA), the experimental 
 setup will be installed on a sounding rocket and will perform under reduced gravity and 
 vacuum conditions. The experiment consists of heated liquid 
 water being let into a chamber connected to a de Laval nozzle, initially filled with 
 low-pressure air. The nozzle lid is opened after a short time period and this 
 depressurisation of the chamber is expected
 to result in vaporisation of the water present in the evaporation chamber.

%We have developed a fluid dynamics model based on Schmidt et al. (2008), to describe the plumes
%characteristics along icy crevasses that could exist on Enceladus. This model accounts for particle
%production generated from homogeneous nucleation, particle growth, wall accretion and sublimation
%and the viscous interaction with the walls. The wall interactions, reservoir conditions and the
%geometry of the channel are studied to reproduce some characteristics of the plumes observed by
%Cassini. Icy crevasses’ geometry will evolve with time due sublimation of ices from the walls or
%accretion onto the walls. The evolution of the geometry of the crevasses and the effect on the
%plumes’ characteristics has been determined with our model. In parallel, we developed an
%experimental setup to study the effect of the nozzle geometries on the plumes characteristics and to
%verify experimentally that under the icy moon’s conditions (pressure, temperature) a plume reaches
%hypersonic/supersonic velocities.
\subsection{The Problem} 
\begin{wrapfigure}{r}{0.40\textwidth}
  \centering
  \caption{Schematic of the flow system assembly.}
  \includegraphics[width = 1\linewidth]{./Chapters/Introduction/figs/experiment}
  \label{fig:experiment}
\end{wrapfigure}
The 'flow system assembly' in the experimental setup, as shown in~\fig{fig:experiment}, 
 consists of the hydraulic accumulator, 
 the evaporation chamber, and the injection 
 system subassembly that drives the liquid from the hydraulic 
 accumulator into the evaporation chamber. In the current 
 design of the setup, these components are assembled vertically and the 
 water is driven up into the evaporation chamber. The evaporation chamber is constructed 
 out of aluminium and can 
 be assumed to be separated into a water reservior and the nozzle.~\COMMENT{-dimensions of reservoir, nozzle. CAD in appendix; where are flow measurements made}

The hydraulic accumulator and the evaporation chamber can be heated until lift-off. Current design considerations 
 assume the temperature of the evaporation chamber and water in the accumulator to be $368.15$\,\si{K} and $353.15$\,\si{K} respectively. 
 The ambient temperature at lift-off is assumed to be $280$\,\si{K} based on the remarks and considerations made during reviews of one such previous FH Aachen experiment.
 The flow system assembly is sealed off before lift-off and an intial pressure of $1$\,\si{bar} is assumed in the evaporation chamber, initially 
 occupied by air. The hydraulic accumulator is sufficiently insulated and the nozzle is encased in a golden foil to minimise the amount of 
 heat lost. No stipulations have been made regarding the evaporation chamber, however as there is no fluid convection
 during operation, losses are assumed to be minimum. 


Water injection begins after the electronics module signals the start of data storage and ends shortly before the apogee is reached. 
 Due to water remaining at the time of cutoff in the injection subassembly and also the accumulator, the total amount of water discharged to 
 the chamber is assumed to be $150$\,\si{mL}. A constant mass flow rate of $15$\,\si{mL/s} is assumed. The inflow velocity in the actual experiment can be tuned to 
 desired values. Discharge time in the simulation is calculated based on the estimated water level at the end of discharge.
 A minimum time of $6$\,\si{s} and a system pressure
 of $1.7$\,\si{bar} is assumed before the electronics system signals the opening of the nozzle lid ensuing the vaporisation of the water.
 As no active heating is provided, the phase change also takes place in the fluid bulk. This process is facilitated by the presence 
 of nucleation points that are found in the water as $2.5$\,\si{g} of silicates per $100$\,\si{mL} of water. The silicates are assumed to be spherical 
 with a uniform diameter of $0.8\,\mu$\si{m}. At the start of depressurisation, these 'seed' points are assumed to be uniformly
 distributed in the water bulk.

%The physical quantities listed above - velocity, pressure, and temperature - consist of the main unknowns 
% that drive the flow characteristics of the system. 
% Characteristically similar to a free surface problem, the first subproblem is modelled 
% using the volume-of-fluid method. 
 %The first subproblem is modelled as a two-phase problem and uses the volume-of-fluid method. The 
 %Considerations about the gravitational force are based on the accelerometer data
 %recorded during previous REXUS flights and are made .~\citelater{newSED} T%he phases involved 
%he evaporation chamber is constructed out of Aluminium-6061 has a stipulated height, width, and depth of $128$, $72.34$, and $50$\,mm respectively.
 %A technical drawing of the evaporation chamber is attached in the appendix~\COMMENT{fake news}. 
%No flow measurements are made in the injection assembly. Pressure and temperature are measured in the evaporation chamber along 
% the nozzle at the midspan, the throat, and the exit.
%measurements in the simualation are made using section probes instead of point probes because section probes are essentially infinite point probes in 2D

\subsection{Multiphase Flows and Star-CCM+}
Multiphase flows are characterised by the simultaneous flow of several phases. Regardless of our 
 improved understanding 
 of such flows, multiphase flows remain complicated to model. Multiphase flows are 
 multiscale in nature - where 
 the cascasding effects of various physics acting over 
 system-level to the microscopic scale must be taken into account~\cite{ishiiThermoFluidDynamicsTwoPhase2011}. 
 Along with turbulence, multiple deformable, moving interfaces pose discontinuities 
 in the fluid properties and induce fluctuations in the flow variables in the neighbourhood. 
 The distribution of the phases cannot be accurately estimated
 a priori~\cite{yadigarogluNatureMultiphaseFlows2018} but the macroscopic 
 behaviour of the system is dependent on small-scale interactions of the 
 phases~\cite{hanrattyPhysicsGasLiquid}. These evolving interfacial structures tend
 to directly affect the inertial and thermodynamical development of the flow field.~\cite{ishiiThermoFluidDynamicsTwoPhase2011}

Two-phase flow can be classified into three main flow regimes based on the 
 distribution of the phases - dispersed flow, transitional flow, and separated flow.
 In case of liquid-gas flows, this ranges from bubbly flows - where small gas bubbles are finely dispersed through the
 continuous liquid phase, to the separated flows - where a sharp interface
 marks the boundary between the two phases.
 The transitional flow regime is marked by the presence of both dispersed and 
 separated flows.~\cite{yadigarogluNatureMultiphaseFlows2018}

Star-CCM+ is a commerical CFD software currently maintained and developed 
 by Siemens Digital Industries Software. The software uses the Finite Volume Method to
 discretise the integral form of the governing differential equations over a number 
 of control volumes in the domain. The software has an automatic meshing system but allows
 the implementation of custom meshing using surface wrappers.~\cite{star-ccm+}
 \COMMENT{add small para about different physics and modelling capabilities.}

\subsection{Objectives and Outline}
~\COMMENT{--------under construction ---------------}\hfill\\
\begin{itemize}
    \item we describe the theoretical framework involved in modelling the 
     two phenomenon
    \item we simulate the 1D stefan problem and 2D RT instability
  using the vof method. we also compare the emp framework in its ability to 
  simulate the 2D RT instability? 
    \item we simulate the infiltration of the evaporation chamber using the 
    vof method,
    discuss the spatial and temporal convergence and a qualitative
    behaviour of the flow field. 
\end{itemize}
The two main physical processes in the experiment - infiltration and evaporation - not only occur at distinct
 points in time but also require distinct mathematical treatment in their simulation. Thus, dividing 
 the experiment into two subproblems -  
 \begin{itemize}
    \item the infiltration of heated water in the evaporation chamber exposed to ambient temperature conditions, 
    \item and the evaporation of superheated water in the evaporation subject to depressurisation - 
 \end{itemize}
  this thesis investigates the feasibility of modelling the two processes in the commercial CFD software tool - Star-CCM+.






